\section{The all-spin background and the constrained linear packet}
\label{sec:background}
\label{bg:section}

Fix $d\in\mathcal U$ and let $g$ be its smooth vacuum metric on the
comparison development of Proposition~\ref{test:comparison}, at the
normalized center $0<\mathfrak a<1$. Given the branch selected in
Section~\ref{test:section}, our objective is to construct two versions
of one linear perturbation: an oscillatory tensor near a late branch
point and an exterior beam whose data lie on the original bridge.
Their full spacetime jets must agree to prescribed finite accuracy at
the entry cylinder. The next
section will turn the two signs of these approximate fields into exact
vacuum data.

We first record the background estimates imported from
\cite{Foundation}, then locate the branch observation and the rays
returning its entry data to the bridge. Only then do we introduce the
normalized tensor slots used in the wave equation and construct the
packet. The comparison construction fixes its inner profile depth
$q_0$ before the exterior computing cuts and the $p_{10}$ radius.
We retain those choices throughout; the return-ray proof below recalls
the depth condition that ensures escape. All these background choices
are independent of the finite packet parameters.

\begin{definition}[Fixed-order subexponential bounds]\label{bg:subexponential}
The notation $Q_T\le e^{o(T)}$ means that, for every $\sigma>0$,
$Q_T\le C_\sigma e^{\sigma T}$ for all sufficiently large $T$.
The constants may depend on $d$, the fixed coordinate domains, and
every fixed derivative and formal expansion order.  A positive
quantity has two-sided subexponential size if this holds for the
quantity and its reciprocal.  A finite product or sum of such bounds,
including a factor polynomial in $T$, is again subexponential.
No uniformity over unbounded derivative orders is intended.
\end{definition}

\subsection{Profiles, entry data, and the full double-null wedge}

Write the interior metric and optical variables as
\begin{equation}\label{bg:wedge}
\begin{gathered}
 t=(u+v)/2\ge0,\qquad x=(v-u)/2,\qquad \mathcal S=\{t=0\},\\
 g=-2a\,du\,dv+
 \gamma_{AB}(d\theta^A-b^Adu)(d\theta^B-b^Bdu),\qquad
 a=q\mathcal A,\quad q=q_0e^{-2\kappa t},\\
 L_0=\partial_u+b^A\partial_A,\qquad L_1=\partial_v.
\end{gathered}
\end{equation}
The constants $q_0,\kappa>0$ depend on the fixed center; the optical
coefficient $a$ is distinct from the Kerr rotation $\mathfrak a$.
Use a fixed finite smooth sphere atlas.  If $y_0=u$, $y_1=v$,
$B_0=b$, and $B_1=0$, put
\begin{equation}\label{bg:null-fields}
\begin{gathered}
 \mathfrak L_i=\partial_{y_i}+\mathcal L_{B_i},\qquad
 \chi_i=\tfrac12\mathfrak L_i\gamma,\qquad
 \tau_i=\operatorname{tr}_\gamma\chi_i,\qquad
 \ell_i=L_i\log a,\\
 \xi_0+\xi_1=d_\theta\log a,\qquad
 \xi_0-\xi_1=-a^{-1}\gamma(\partial_vb,\cdot),\qquad
 \beta_i=\operatorname{div}_\gamma\chi_i-d_\theta\tau_i.
\end{gathered}
\end{equation}
Angular Sobolev norms are taken in that atlas; the controlled angular
metrics give equivalent geometric norms.  Ordinary derivatives $D$
mean coordinate derivatives of the displayed arrays.

\begin{proposition}[All-spin geometric interface]\label{bg:geometry}
There is a choice of $\epsilon>0$, depending only on the fixed Kerr
center and the fixed preparation, such that every
$p_{10}(d-d_*)<\epsilon$ in the stated smooth symbol class has the
following properties.
\begin{enumerate}
\item The cylinder $\mathcal S$ in Proposition~\ref{test:comparison}
has a smooth vacuum development on the full product wedge
$[0,\infty)_t\times\mathbb R_x\times\mathbb S^2$ in
Equation~\eqref{bg:wedge}.  The quantities
$\mathcal A^{\pm1},\gamma^{\pm1}$ are uniformly bounded in angular
$H^6$, and $\xi_i,a^{-1}\partial_vb$ in angular $H^5$.
There are nonnegative functions $g_i(y_i)$ with summable unit-cell
suprema such that, in the indicated low angular norms,
\begin{equation}\label{bg:envelopes}
 \|\chi_i,\ell_i+\kappa\|_{H^4}\le C(q+g_i(y_i)),
 \qquad \|b\|_{H^5}\le C(q+g_0(u)).
\end{equation}
\item The following quantities are uniformly bounded, and tend to
zero when $t,|u|,|v|\to\infty$:
\begin{equation}\label{bg:low-limits}
\begin{gathered}
 \|b\|_{H^5}+\sum_i\|\chi_i,\ell_i+\kappa\|_{H^4},\\
 \sum_{k=0}^1\left(\|\partial_{y_k}b\|_{H^2}
       +\sum_i\|\partial_{y_k}\chi_i,
                          \partial_{y_k}\ell_i\|_{H^2}\right),\\
 \sum_{k=0}^1\left(\|\partial_{y_k}\gamma\|_{H^4}
                 +\|\partial_{y_k}\log a+\kappa\|_{H^4}\right).
\end{gathered}
\end{equation}
In particular $D\log a,D\gamma,b,Db$ are uniformly bounded.
\item For each fixed $U\in\mathbb R$, $C>0$, and finite order,
all ordinary derivatives of
\begin{equation}\label{bg:high-arrays}
 \gamma^{\pm1},\quad\log\mathcal A,\quad b,
                  \quad a^{-1}\partial_vb
\end{equation}
are bounded by $e^{o(T)}$ on
\begin{equation}\label{bg:truncation}
 t\ge0,\qquad y_d\le U,\qquad y_l\le CT,
 \qquad \{d,l\}=\{0,1\}.
\end{equation}
This includes every restriction and coordinate-path subsegment.
Every fixed derivative of $a$, $\sqrt a$, or $\partial_vb$ retains,
respectively, the factor $a$, $\sqrt a$, or $a$.
\item Either branch can be put in the convention
$v\to\infty$, $u\to u_*<\infty$.  Exchanging the optical labels
and preparing the angular coordinates anew preserves all these
conclusions on the same product spacetime and initial cylinder.
\end{enumerate}
Only $p_{10}$ is required small; no higher seminorm, radiation
coefficient, or curvature measured by a timelike observer is required
small or bounded.
\end{proposition}

\begin{proof}
\emph{Exterior evolution and cylinder data.}
The exterior input is the simultaneous closure in
\cite[\CompanionRef{foundation}{far:closure}]{Foundation}.  It applies
to both computing cylinders from the $p_{10}$ input, with variable
asymptotic parameters and arbitrary finite higher symbol seminorms.
The higher estimates are
\cite[\CompanionRef{foundation}{co2:higher}]{Foundation}.
We describe the entry hypotheses so that the subsequent interior
citations have a precise common domain of applicability.

For nearby physical Kerr parameters the exact profiles in
\cite[\CompanionRef{foundation}{in:profiles}]{Foundation} are smooth in
$(q,\theta)$ and the parameters, with $\mathcal A,\gamma$ and their
inverses smooth at $q=0$.  They satisfy
\[
 b=O(q),\quad \chi_i,\ell_i+\kappa,\beta_i=O(q),
 \qquad \xi_i,K_\gamma,\operatorname{curl}_\gamma\xi_i=O(1),
\]
after each fixed angular or parameter derivative; here $K_\gamma$ is
the Gaussian curvature of the angular metric. The common $\kappa$ is obtained by rescaling the exact optical variables by
$\kappa_-/\kappa$, where $\kappa_-$ is the positive absolute inner
surface gravity of the current physical profile.

The moving preparation uses integrated Killing and rotation phases,
whose derivatives are the instantaneous frozen frequencies.  Its
Jacobian errors contain positive parameter derivatives with bounded
coefficients, without an elapsed-time factor.  This is
\cite[\CompanionRef{foundation}{in:moving-coordinates}]{Foundation}.
The single physical profile on the wedge depends only on $u$;
a frozen differentiated profile means differentiation at fixed
parameters followed by evaluation at this profile.  The additional
parameter defect is consequently confined to the $u$ direction.
After decreasing $\epsilon$, the cylinder differences from these
frozen profiles have arbitrarily small sums of unit-$x$-cell suprema
in exactly the following norms:
\begin{equation}\label{bg:entry-orders}
\begin{array}{c|c}
 \mathcal A,\log\mathcal A,\gamma,b,\chi_i,\ell_i,\xi_i&H^6\text{ in angle}\\
 \gamma&H^7\text{ in angle}\\
 \partial_x(\chi_i,\ell_i,\xi_i,b)|_{\mathcal S}&H^2\text{ in angle}.
\end{array}
\end{equation}
The unit-cell sum of the first profile derivative is small as well.
Every fixed higher entry norm on $|x|\le CT$ is subexponential.
These are the full assumptions and conclusions of
\cite[\CompanionRef{foundation}{in:entry}]{Foundation}.
Its trace estimate uses mixed spacetime metric derivatives through
order eight and one additional time derivative outside the trapping
region, supplied by \CompanionRef{foundation}{co2:sliced} there.  The entry supports lie between the
horizons, separated from that trapping region.  On the bounded middle
the finite bridge evolution supplies the same derivative budget.
Thus the trace step does not require a small high-order norm beyond
the stated $p_{10}$ input.

\emph{Interior continuation and derivative bounds.}
\CompanionRef{foundation}{in:low}, \CompanionRef{foundation}{in:pure-jets}, and \CompanionRef{foundation}{in:higher} of \cite{Foundation} apply
to precisely these profiles and entry norms.  They give
Equations~\eqref{bg:envelopes}--\eqref{bg:truncation}.  In particular
the second line of Equation~\eqref{bg:low-limits} includes the pure
longitudinal derivatives $\mathfrak L_i\chi_i$ and $L_i\ell_i$;
one has not inferred them by differentiating an uncontrolled limit.
\CompanionRef{foundation}{in:continuation} first uses finite-diamond
high estimates, where $a$ is bounded below, to continue and exhaust
the full wedge.  Only after that exhaustion does it use the tails of
the summable low envelopes to prove the expanding-domain
subexponential estimate.  This order avoids assuming a global tail
estimate to prove existence of the region carrying that tail.
The factor assertions follow by differentiating $a=q\mathcal A$ and
$\partial_vb=a\gamma^{-1}(\xi_1-\xi_0)$.

\emph{The other branch convention.}
Finally \cite[\CompanionRef{foundation}{in:label-swap}]{Foundation} negates
the prepared $x$ on $\mathcal S$ and transports the newly prepared
angles along the new shift-free direction.  At each finite point
this is a finite smooth sphere flow with a smooth inverse.
The entry proof and finite-diamond estimates are reapplied in that
gauge and exhausted.  No uniform bound on an otherwise unestimated
long angular coordinate change is used. These are estimates in the
newly prepared coordinates on the same comparison spacetime.
\end{proof}

\subsection{An arbitrary branch observer}

The selected branch determines the observation point and its optical
scale. The following argument uses only the low geometric estimates;
in particular, selecting the observer requires no curvature assumption.

\begin{lemma}[Branch normalization without a curvature hypothesis]
\label{bg:branch}
Let a future unit timelike geodesic approach a branch
$v\to\infty$, $u\to u_*<\infty$, in the comparison wedge.
With proper time denoted by $\tau$, put
\[
 p=a\dot v,\qquad \Xi=\dot\theta-b\dot u,\qquad P=\gamma\Xi.
\]
Then $p\asymp1$, and $P$, $\Xi$, and $\dot u$ stay bounded on
a late segment.  The angular position has a limit.  At the point
$p_T$ with $v=T$ one has
\begin{equation}\label{bg:branch-scales}
             a_p:=a(p_T)\asymp e^{-\kappa T},
             \qquad \dot v(p_T)\asymp e^{\kappa T}.
\end{equation}
\end{lemma}

\begin{proof}
This is the observer-independent calculation in
\cite[\CompanionRef{foundation}{pk2:branch-normalization}]{Foundation}.
Unit normalization gives
\begin{equation}\label{bg:unit-branch}
        1+|\Xi|_\gamma^2=2a\dot u\dot v=2p\dot u.
\end{equation}
For the unit tangent $U$, differentiating $p=-g(U,L_0)$ and
using the geodesic equation gives the exact identity
\[
 \frac{\dot p}{p}=\ell_0\dot u+(\xi_0-\xi_1)(\Xi)
                           -\frac{\chi_0(\Xi,\Xi)}p.
\]
The null connection identities in Equation~\eqref{test:null-connections} therefore imply
$|\dot{\log p}|\le C(\dot u+|\Xi|_\gamma)$, using
$|\Xi|^2/p\le2\dot u$ from Equation~\eqref{bg:unit-branch}.
Moreover
\[
 \int|\Xi|_\gamma\,d\tau
 \le\sqrt2\left(\int du\right)^{1/2}
                  \left(\int a\,dv\right)^{1/2}<\infty.
\]
Indeed $u$ has finite increasing variation and
$a\le Ce^{-\kappa v}$ on the bounded-$u$ segment.  Integration
therefore bounds $p$ above and below by positive constants.
Since $b$ is bounded, this also proves finite angular length and
an angular limit.

Work near that limit in a sphere chart and introduce canonical
momenta $\pi_u=-p-b\cdot P$.  The unit shell solved with $v$ as
time has Hamiltonian
\begin{equation}\label{bg:branch-hamiltonian}
 -\pi_v=\frac{a(1+|P|_{\gamma^{-1}}^2)}{2p},
 \qquad \frac{du}{dv}
       =\frac{a(1+|P|_{\gamma^{-1}}^2)}{2p^2}.
\end{equation}
At fixed canonical $(\pi_u,P)$, $\partial_\theta p
=-(\partial_\theta b)\cdot P$.  Bounded first derivatives of
$\log a,\gamma,b$ in Hamilton's equation for $P$ consequently give
\[
 \left|\frac{dP}{dv}\right|
        \le C(p+|P|)\frac{du}{dv}.
\]
Gronwall in the finite variation variable $u$ bounds $P$.
Equation~\eqref{bg:unit-branch} bounds $\Xi$ and $\dot u$.
Finally Equation~\eqref{bg:wedge} yields
Equation~\eqref{bg:branch-scales}.  This argument uses only first
metric derivatives and the unit shell, and no curvature bound along
the geodesic.
\end{proof}

\subsection{The exterior geometry of the launch patch}

\begin{proposition}[Entry timing, escape, and flow estimates]
\label{bg:exterior}
On either prepared end use regular exterior time $t_e$, co-moving
coordinates $y$, and lab coordinates $z=y+C_e(t_e)$, with fixed
$\sup|C_e'|\le s_*<1/10$, after decreasing $\epsilon$ if needed.  The relevant slice $t_e=0$ agrees with the
original bridge sufficiently far out.  There is a preparation
parameter $I$ with $0<c_I\le I'\le C_I$ and $I(t_T)=T$.
The launch patch near $v=x=T$ on $\mathcal S$ occurs at
$t_e=t_T+O(1)$, and $t_T\asymp T$.  A fixed enlargement is
contained in $t_e\le t_T+C_2$.  Its selected-end portion
$v<(1-\nu_0)T$ is earlier by at least $c\nu_0T-O(1)$ for each
fixed $\nu_0>0$.

The exterior metric has uniformly regular spacelike temporal
levels.  Its first and second ordinary lab derivatives are bounded
on the computing region and become arbitrarily small at sufficiently
large co-radius, uniformly in time.  Every fixed mixed jet on the
used time and radius ranges $O(1+T)$ is $e^{o(T)}$.  Entry charts
have a noncharacteristic collar with controlled low jets and
subexponential fixed higher jets; the inner exit has a strict margin
deeper than the entry cylinder.

The null covectors $dv$ launched on a sufficiently small fixed patch
at this entry escape towards the past to the same original end.
For each fixed $\beta>0$, after fixing a sufficiently large radius
and taking $T$ large, their bridge arrival points satisfy
\begin{equation}\label{bg:arrival-shell}
 (1-\beta)t_T<|z-z_*|<(1+\beta)t_T,
                  \qquad z_*=C_e(t_T).
\end{equation}
The entry crossing is unique and transverse, and separated from
other used ray portions outside a uniform collar.  Every fixed
Hamiltonian flow jet, including launch-center derivatives and the
propagators on all subsegments, is $e^{o(T)}$.  Tensor parallel
transport and its inverse are bounded.  These assertions hold on
short past and future launch enlargements through the deeper exit.
\end{proposition}

\begin{proof}
The timing and low exterior bounds are established in
\cite[\CompanionRef{foundation}{sec:compact}, \CompanionRef{foundation}{sec:far}, and \CompanionRef{foundation}{pk2:return-rays}]{Foundation}; the escape and all-subsegment
statements are \CompanionRef{foundation}{pk2:escape} there.
We check why the launch geometry meets that proposition independently
of the timelike observer. The original profile-depth choice makes the
radial potential of these particular null rays strictly positive.
For one of the nearby exact Kerr profiles, write $M_K$ and
$\mathfrak a_K$ for its mass and spin length, $r_K$ for the
Boyer--Lindquist radius, $\omega=(\vartheta,\varphi)$ for the
angular variables, and $|\cdot|_{\mathrm{rd}}$ for the unit round
sphere norm. Thus
\[
 \Delta_K=r_K^2-2M_Kr_K+\mathfrak a_K^2,\qquad
 r_\pm=M_K\pm\sqrt{M_K^2-\mathfrak a_K^2}.
\]
In aligned Kerr coordinates its optical correction $F_K$ solves
\begin{equation}\label{bg:profile-eikonal}
 2(r_K^2+\mathfrak a_K^2)\partial_{r_K}F_K
 +\Delta_K(\partial_{r_K}F_K)^2
 +|d_\omega F_K|_{\mathrm{rd}}^2
 +\mathfrak a_K^2\sin^2\vartheta=0,
 \qquad F_K(r_-)=0.
\end{equation}
The finite root is noncharacteristic.  In the original comparison
preparation, $q_0$ is chosen small enough for both angular transport
conventions and for the uniform strict depth condition
\begin{equation}\label{bg:profile-depth}
 \left(|\partial_\vartheta F_K|^2
              +\mathfrak a_K^2\sin^2\vartheta\right)_{\rm entry}
                   <4\mathfrak a_K^2.
\end{equation}
This is possible because the physical spins lie in a compact
neighborhood inside $(0,1)$ and $\partial_\vartheta F_K=O(q_0)$.
The inner exit is fixed strictly deeper, with a positive margin
before the inner horizon.

On fixed radial compact sets the
metric and prepared entry covectors converge in the required low
orders to a nearby frozen subextremal Kerr metric and to
$\lambda_-d(r_*+F_K+T_K)$ in aligned variables, with
$\lambda_-=\kappa_-/\kappa>0$. Here $T_K$ is the signed
Boyer--Lindquist Killing time and the tortoise coordinate $r_*$
satisfies $dr_*/dr_K=(r_K^2+\mathfrak a_K^2)/\Delta_K$.
The integrated phase preparation used above gives
this covector convergence even if accumulated azimuthal phases do
not converge.  The frozen axial momentum is zero and its radial
potential, divided by the square of Killing energy, is
\begin{equation}\label{bg:radial-potential}
 (r_K^2+\mathfrak a_K^2)^2-\Delta_K
     \left(|\partial_\vartheta F_K|^2
                 +\mathfrak a_K^2\sin^2\vartheta\right)_{\rm entry}.
\end{equation}
For $r_K\ge r_+$, $\Delta_K<r_K^2$ and
$(r_K^2+\mathfrak a_K^2)^2/r_K^2\ge4\mathfrak a_K^2$;
between the horizons $\Delta_K<0$.  The prior depth choice
Equation~\eqref{bg:profile-depth} therefore makes this potential
strictly positive.  Its initial past radial direction is outward,
since the finite Hamilton--Jacobi root gives
$r_K^2+\mathfrak a_K^2+\Delta_K\partial_{r_K}F_K>0$.
It crosses each fixed compact radial band, including the future
event horizon in regular coordinates, transversely in bounded
regular time.  The defining entry level is strictly monotone near
its crossing; the other used compact portions are separated from it.
These are uniform launch-patch facts, and persist by the compact
convergence just stated.

In the far region the proof of \CompanionRef{foundation}{pk2:escape} gives a summably
small connection integral on the outward-to-the-past cone.  More
precisely, \CompanionRef{foundation}{pk2:far-ray-coefficients} there gives the far-region lab-derivative estimate: it controls
the bad first-derivative term on each radial dyad $r\asymp R$ by
an integrated bound
$O(R^{(e_f-(1-\delta_f)A_H)/2})$, with
$e_f-(1-\delta_f)A_H<0$, and the remaining term is summable.
Here $e_f,\delta_f,A_H$ are the fixed far-estimate exponents of
that source.  Thus the lab velocity stays close to a fixed outward
unit direction and the frequency stays nondegenerate.  Integration
gives Equation~\eqref{bg:arrival-shell}; the same connection
integral bounds tensor transport and its inverse.
For the Hamiltonian first variation, only the momentum-to-position
block need remain of order one at large radius.  Rescaling its
position variation by a small fixed constant makes that block
small, and then taking the far radius large makes the other blocks
small by the first and second metric derivative bounds.  The
complementary compact crossing has fixed length.  This proves an
arbitrarily small average growth rate on every subsegment.
Higher variation equations have this same homogeneous propagator
and sources involving lower jets and fixed high coefficient jets;
induction gives the asserted subexponential bounds.  The proof
uses only the prepared launch covectors and the exterior geometry,
so applies at the moving angular foot of any branch and after
the separately prepared exchange of labels.
\end{proof}

\subsection{Normalized tensor slots and wave coefficients}

We now prepare the tensor wave equation for the packet. In the
coordinate frame, inverse metric coefficients grow as the optical lapse
$a$ decreases. The following frame separates that geometric scale from
the bounded coefficient arrays. It also identifies the constant radial
connection terms that must be retained in the interior energy estimate.

For $e_A=\partial_{\theta^A}$ define
\begin{equation}\label{bg:frames}
 \hat e_i=a^{-1/2}L_i,\qquad \hat e_A=e_A,\qquad
 Z_\alpha=\sqrt a\,\hat e_\alpha,\qquad
 W_t=L_0+L_1,\quad W_x=L_1-L_0.
\end{equation}
Thus $Z_i=L_i$, $Z_A=\sqrt a\,e_A$, and the normalized metric
matrix has radial null block with $g(\hat e_0,\hat e_1)=-1$ and
angular block $\gamma$.  A relative array of a covariant two-tensor
$h$ means $h(\hat e_\alpha,\hat e_\beta)$.  Derivatives of this
array remain ordinary $D$ derivatives.  For a covariant two-tensor
entry $n$, let $n_i$ count its slots of index $i$, and set
\begin{equation}\label{bg:slot-weight}
                c_n=\kappa(n_0-n_1)/2,\qquad |c_n|\le\kappa.
\end{equation}

\begin{lemma}[Normalized coefficient arrays]\label{bg:coefficients}
Let $V_\alpha$ be the matrix of $\nabla_{Z_\alpha}$ on the
normalized vector basis.  These matrices are bounded, and in the
bulk limit $\min(v,-u,t)\to\infty$,
\begin{equation}\label{bg:connection-limits}
 V_0\longrightarrow D_0=\operatorname{diag}(-\kappa/2,\kappa/2,0,0),
 \qquad V_1\longrightarrow-D_0,\qquad V_A\longrightarrow0.
\end{equation}
The differentiated connection terms contracted in the tensor
operator $a\square_g$, and every normalized entry of
$a\operatorname{Riem}(g)$, are bounded and tend to zero in this
limit.  Their fixed ordinary jets and those of the brackets of $Z$
are $e^{o(T)}$ on Equation~\eqref{bg:truncation}.  Moreover
$\nabla_{L_1}L_1=\ell_1L_1$ exactly.
The undifferentiated constant connection products are retained:
on entry $n$ the limiting radial operator is
\begin{equation}\label{bg:radial-conjugation}
\begin{split}
 &-(L_0+c_n)(L_1-c_n)-(L_1-c_n)(L_0+c_n)\\
 &\qquad=-L_0L_1-L_1L_0+2c_n(L_0-L_1)+2c_n^2.
\end{split}
\end{equation}
It is the radial part of
$e^{2c_nx}a\square_g^{\rm scal}e^{-2c_nx}$.
\end{lemma}

\begin{proof}
The needed calculations are those of
\cite[\CompanionRef{foundation}{pk2:coefficients}]{Foundation}; we display
them to identify every possible unsuppressed term.  Koszul's formula
applied to Equations~\eqref{bg:wedge} and \eqref{bg:null-fields}
gives, with $j=1-i$,
\begin{align}
 \nabla_{L_i}\hat e_i&=\tfrac12\ell_i\hat e_i,&
 \nabla_{L_i}\hat e_j&=-\tfrac12\ell_i\hat e_j
                         +\sqrt a\,\xi_i^Ae_A,\label{bg:radial-connection}\\
 \nabla_{L_i}e_A&=(\chi_{i,A}{}^B-\partial_AB_i^B)e_B
                         +\sqrt a\,\xi_{i,A}\hat e_i,\notag\\
 \nabla_{Z_A}\hat e_i&=\chi_{i,A}{}^Be_B
          +\tfrac12\sqrt a(\xi_i-\xi_j)_A\hat e_i,&
 \nabla_{Z_A}e_B&=\sqrt a\,\Gamma(\gamma)_{AB}^Ce_C
             +\chi_{0,AB}\hat e_1+\chi_{1,AB}\hat e_0.\notag
\end{align}
For example, differentiating $L_i=\sqrt a\,\hat e_i$ gives
$\nabla_{L_i}L_i=\ell_iL_i$.  The low estimates yield
Equation~\eqref{bg:connection-limits}.  The exact brackets are
\begin{equation}\label{bg:brackets}
\begin{gathered}
 [L_0,L_1]=\sqrt a(\xi_0-\xi_1)^AZ_A,\\
 [L_i,Z_A]=\tfrac12\ell_iZ_A-(\partial_AB_i^B)Z_B,\\
 [Z_A,Z_B]=\tfrac12\sqrt a\bigl(
       (\partial_A\log a)Z_B-(\partial_B\log a)Z_A\bigr).
\end{gathered}
\end{equation}
The radial trace of differentiated connection coefficients uses
$L_0V_1$ and $L_1V_0$, and hence cross derivatives.  In particular
$L_1\partial_Ab=\partial_A\partial_vb$ retains $a$.
Longitudinal derivatives of torsion reduce by the vacuum identities
\[
 \mathfrak L_i\xi_j=\beta_i+\tau_i\xi_i,
 \qquad \mathfrak L_i\xi_i=d_\theta\ell_i-\beta_i-\tau_i\xi_i.
\]
The remaining angular derivatives carry $\sqrt a$.
These observations and Equation~\eqref{bg:low-limits} give the
boundedness and bulk smallness of the contracted differentiated
connection terms.

Writing $V_X$ for the connection matrix along $X$, curvature obeys
the exact matrix identity
\begin{equation}\label{bg:curvature-matrix}
 aR(\hat e_\alpha,\hat e_\beta)
 =Z_\alpha V_\beta-Z_\beta V_\alpha
       +[V_\alpha,V_\beta]-V_{[Z_\alpha,Z_\beta]}.
\end{equation}
Besides the cross derivatives just considered, its radial-angular
entries contain the pure derivatives $L_i\chi_i$ controlled by
Proposition~\ref{bg:geometry}.  Two angular directions carry a
$\sqrt a$ derivative or a product of matrices tending to zero.
The only nonzero limiting bracket coefficient multiplies
$V_A\to0$, while $[D_0,-D_0]=0$.  Angular $H^2$ embeds into
$C^0$ on the controlled spheres, so the stated low derivative
budget suffices.  Differentiating these identities and using
Equation~\eqref{bg:high-arrays} proves the higher-jet assertions.
The covariant tensor connection has the negative vector connection
on each slot; its constant radial actions are therefore
$L_0+c_n$ and $L_1-c_n$.  Expanding their mixed trace gives
Equation~\eqref{bg:radial-conjugation}, since
$L_0x=-1/2$ and $L_1x=1/2$.  In particular $2c_n^2$ has not
been classified as a small coefficient.  The angular metric and
rescaled torsion need not tend to a prescribed constant or be small.
\end{proof}

\subsection{A constrained packet with freely chosen amplitude}\label{bg:linear-packet-section}

Fix a branch point $p_T$ as in Lemma~\ref{bg:branch}. The construction
is linear in its amplitude. We state it at the amplitude used in the
present proof and record its general-amplitude consequence at the end
of the section. Choose
$u_2>u_*$, $c_2>0$, and then small $c_0>0$ so that $p_T$ has
a fixed interior margin in the finite domain
\begin{equation}\label{bg:packet-domain}
\begin{gathered}
 \mathcal D_T=\left\{0\le t\le L_T=(T+u_*)/2+c_0,\quad
 u+c_0t/T\le u_2,\quad v+c_0t/T\le T+c_2\right\}.
\end{gathered}
\end{equation}
We use $\mathcal D_T$ as the closed domain on which to state packet
bounds; every construction is first made on a slightly enlarged domain,
so its boundary jets come from a spacetime neighborhood. The nonlinear
evolution will be retained inside the strict upper-time and side
inequalities of the same domain. Its slice intervals in $x$ have
lengths bounded below. Fix $u_f>u_2$
for terminal transports.  The whole strip $-v\le u\le u_f$,
with $v$ near $T$, is contained in Equation~\eqref{bg:truncation},
including its moving initial foot $u=-v$.

For a symmetric covariant tensor define trace reversal by
$\operatorname{tr.rev.}_g h=h-\tfrac12(\operatorname{tr}_g h)g$.
Let $\mathcal P$ be the tensor-wave operator with curvature
potential obtained from the linearized wave-map reduced vacuum
equation after trace reversal, normalized to have principal part
$\square_g$.  Write $\mathcal C=\operatorname{div}_g$ on these
trace-reversed tensors.  On the Ricci-flat background,
\begin{equation}\label{bg:linear-bianchi}
                     \mathcal C\mathcal P=\mathcal P_1\mathcal C,
\end{equation}
where $\mathcal P_1$ is the corresponding one-form wave operator.
This is the linearized contracted Bianchi identity.  The underlying
nonlinear reduction throughout is the local wave-map equation with
fixed target $g$:
\begin{equation}\label{bg:wave-reduction}
\begin{split}
 \operatorname{Ric}(\widetilde g)_{\mu\nu}
       +\nabla^{\widetilde g}_{(\mu}\Upsilon_{\nu)}&=0,\\
 \Upsilon_\mu&=-\widetilde g_{\mu\lambda}
 \widetilde g^{\alpha\beta}
 \bigl(\Gamma(\widetilde g)^\lambda_{\alpha\beta}
                         -\Gamma(g)^\lambda_{\alpha\beta}\bigr).
\end{split}
\end{equation}
This is the local wave-map reduction of \cite[\CompanionRef{foundation}{pk2:local-reduction}]{Foundation}.
Its target is a fixed coefficient field on each computing block.

\begin{proposition}[Linear packet and exterior beam]\label{bg:packet}
Fix any finite list of output derivative orders, including an
exterior base buffer $N_b$.  For every fixed $\zeta>0$ and every
subsequently prescribed inverse accuracy $M_1$, sufficiently high
finite transport and Taylor orders yield a real interior tensor
$h_T^p$ and a real exterior tensor $h_T^B$ with the following
properties.  Put
\begin{equation}\label{bg:packet-scales}
       A_T=e^{-\alpha T},\qquad \alpha=7\kappa/8,
       \qquad\lambda_T=e^{\zeta T}.
\end{equation}
The interior field has the form
\begin{equation}\label{bg:packet-form}
 h_T^p=A_T\operatorname{tr.rev.}_g\operatorname{Re}
       \left(e^{i\lambda_Tv}\sum_{j=0}^{J}\lambda_T^{-j}p_j\right).
\end{equation}
Every fixed normalized coefficient jet of $p_j$ is $e^{o(T)}$.
The field is supported in a fixed-width neighborhood of $v=T$,
with angular support in the transported launch patch.  Its relative
array $q_T=(h_T^p(\hat e_\alpha,\hat e_\beta))$ satisfies
\begin{equation}\label{bg:packet-size}
              \|q_T\|_{C^N(\mathcal D_T)}
                    \le A_T\lambda_T^N e^{o(T)}
\end{equation}
at each specified order.  The tensor $p_0$ is trace-free and
$p_0(L_0,\cdot)=0$.  There is a purely angular $\gamma$-unit vector
$X$ at $p_T$ for which
\begin{equation}\label{bg:polarization}
 \sigma_T:=\left|\operatorname{Re}
        \bigl(e^{i\lambda_TT}p_0(X,X)\bigr)\right|
\end{equation}
has two-sided subexponential size.  The $a$-scaled linear wave
residual and $\sqrt a$-scaled linear gauge residual, in normalized
slots and with all prescribed ordinary derivatives, are bounded
by $A_T\lambda_T^{-M_1}e^{o(T)}$.

In regular exterior lab components,
\begin{equation}\label{bg:beam-estimates}
 \|h_T^B\|_{C^{N_b}}\le A_T\lambda_T^{N_b+3/2}e^{o(T)}.
\end{equation}
Its linear wave and gauge residuals, and its difference from
$h_T^p$ on the entry collar, have size
$A_T\lambda_T^{-M_1}e^{o(T)}$ at all the prescribed finite orders.
The latter assertion includes normal spacetime jets.  On the
original bridge its support lies in Equation~\eqref{bg:arrival-shell}
with arbitrarily thin tube margins.  Increasing $M_1$ and the
finite formal degrees does not increase the positive frequency
powers in Equations~\eqref{bg:packet-size} and
\eqref{bg:beam-estimates}.  Thresholds and subexponential constants
may depend on those degrees.
\end{proposition}

\begin{proof}
We extract the linear construction of
\cite[\CompanionRef{foundation}{pk2:deep-packet}, and the proof of \CompanionRef{foundation}{pk2:beam-proposition}]{Foundation}.
The verification below records its amplitude independence and its
uniformity at the actual frequency $\lambda_T$.

\emph{Constrained transport in the interior.}
Set $k_v=g^{-1}dv=-a^{-1}L_0$, $K_vp=p(k_v,\cdot)$, and
$\mathcal T=2\nabla_{k_v}+\square_gv$.  Since $v$ is optical,
the transport and divergence equations are
\begin{equation}\label{bg:hierarchy}
 i\mathcal Tp_j+\mathcal Pp_{j-1}=0,
 \qquad iK_vp_j+\mathcal Cp_{j-1}=0,
 \qquad p_{-1}=0.
\end{equation}
Multiplication of the transport by $-a/2$ gives
$\nabla_{L_0}+\tfrac12\tau_0$.  On a normalized entry of weight
$c_n$ its same-weight constant part is $L_0+c_n$, and the
remaining same-weight coefficients have arbitrarily small bulk
average.  The connection formulas show that off-weight transport
couples only from higher weights.  At $u=-v$ prescribe the slots
with a $0$ index using the gauge equation.  The free angular
entries vanish for $j>0$; for $j=0$ take a compact smooth angular
and $v$ bump times a trace-free polarization of unit size at the
foot of the ray from $p_T$.  Initialize the negative-weight slots
at $u_f$ with zero terminal values and integrate them backwards.
The constants $c_n$ have the damping sign in these respective
directions.  The decreasing-weight triangular system therefore
determines all coefficients.

Conjugate Equation~\eqref{bg:linear-bianchi} by the oscillatory
phase and compare successive powers.  Once preceding gauge
defects vanish, the next defect in Equation~\eqref{bg:hierarchy}
obeys homogeneous one-form transport and has zero initial value.
Hence every gauge equation holds throughout the strip.  The
terminal choices impose no further gauge condition: contraction
by $k_v$ uses exactly the entries with a $0$ slot.

For completeness, along each $L_0$ transport the portion outside
$\min(v,-u,t)>R$ has length bounded in terms of $R,u_f$ only.
On its complement $b$, its first angular derivatives, and the
nonconstant diagonal coefficients become arbitrarily small.
Their integrals are consequently $o(T)$.  The angular flow and
its inverse have subexponential first jets.  Every higher flow
variation solves the same homogeneous equation with lower-jet
sources and the high coefficient bounds of
Equation~\eqref{bg:truncation}; induction gives subexponential
jets, including derivatives of the moving foot.  The favorable
constant signs just noted give the identical conclusion for
each transport coefficient.  The normalized sources
$a\mathcal Pp_{j-1}$ and $\sqrt a\,\mathcal Cp_{j-1}$ have
these bounds by Lemma~\ref{bg:coefficients}.  Integration over
length $O(T)$ preserves them.  All these equations are local in
$v$ and along the angular flow, so preserve the asserted support,
also in the terminal integrations.

The leading tensor remains trace-free and annihilates $k_v$.
Its screen class transports by an isometry of the positive screen
metric times an area factor with logarithmic derivative of
integral $o(T)$.  Its norm thus has two-sided subexponential
size.  A trace-free symmetric form on this two-dimensional screen
has a unit eigenvector with a diagonal value comparable to its
norm.  Represent it by a purely angular $X$.  Adding a multiple
of $L_0$ would not change $p_0(X,X)$, so no observer-dependent
screen adjustment is required.  Multiplication of every $p_j$
by one constant unit complex number makes the diagonal contraction
real at $v=T$, proving Equation~\eqref{bg:polarization}.
Truncation of Equation~\eqref{bg:hierarchy} leaves only the last
transport and gauge remainders.  Taking $J$ large relative to
$M_1$ and the fixed differentiated output orders proves the
interior residual claim.  Each ordinary differentiation has cost
at most one positive power of $\lambda_T$; the factors $a$ and
$\sqrt a$ already account for the normalized wave and divergence
operators.

We have obtained the interior packet, including its polarization,
support, and linear residual bounds. It remains to realize its entry
jets by an exterior field that returns to the original bridge.

\emph{Exterior complex phases and amplitudes.}
Complex Gaussian beam superposition follows the method developed in
\cite{Ralston1982,Tanushev2008,LiuRunborgTanushev2013}. We retain the
phase-space construction and finite matching below because they must
also preserve the tensor gauge and give subexponential bounds on the
growing time interval. Parameterize centers $Y$ by the three
coordinates $w=(x,\theta)$ of the launch patch.  Start a complex
phase $\varphi$ with spatial value $v(0,w)+i|w-Y|^2/2$ and choose its
temporal root to agree with $dv$ at the center.  Proposition~
\ref{bg:exterior} supplies the real rays and every subsegment
propagator.  Propagate finite complex phase jets along them.  If
$P_c,Q_c$ are the position and momentum differentials of this
complex gradient graph, symplecticity gives
\begin{equation}\label{bg:complex-symplectic}
 P_c^*Q_c-Q_c^*P_c=2i\operatorname{Im}M_c,
 \qquad
 \operatorname{Im}(Q_cP_c^{-1})
           =P_c^{-*}(\operatorname{Im}M_c)P_c^{-1}>0,
\end{equation}
where $M_c$ is the initial Hessian.  The initial positive imaginary
part and its inverse are controlled on the noncharacteristic
collar.  The first identity bounds $P_c^{-1}$ by the bound for
$Q_c$, and the second bounds the inverse of the propagated
imaginary Hessian by the bound for $P_c$ squared, times the
initial inverse.  Thus all these quantities are $e^{o(T)}$, even
when the real spatial projection has caustics.  Finite Taylor
composition and reversion use only smooth jets at real points.

Central amplitudes transport by tensor parallel transport times
the inverse square root of the density ratio determined by
$\sqrt{|\det g|}\,k^{t_e}\det P_c$, where $k=g^{-1}d\varphi$
is evaluated at the beam center. This propagator and its
inverse are subexponential on every subsegment; higher jets have
the same homogeneous propagator and lower-jet sources.  Impose
both wave and gauge transports in a sufficiently high finite
Taylor quotient.  The conjugated Bianchi identity is their defect
compatibility identity.  Contraction $Kp=p(k,\cdot)$ has the
explicit right inverse
\begin{equation}\label{bg:contraction-inverse}
 R_kf=l\otimes f+f\otimes l-f(k)l\otimes l,
                         \qquad l(k)=1.
\end{equation}
The nonzero pairing remains invertible in formal jets.  This
fixes the forced gauge jets and leaves the projected kernel jets
free.

The complex phases now remain nondegenerate through real caustics,
and the constrained amplitude transport has identified the free kernel
jets. We use those jets to match the interior field at entry.

\emph{Matching the full spacetime jets.}
Superpose the beams with center measure
$(\lambda/2\pi)^{3/2}\,dY$.  On the unique entry collar the
center projection is a diffeomorphism.  Changing the center
variable to its position and taking Gaussian moments produces
$e^{i\lambda v}\sum_j\lambda^{-j}\mathfrak p_j$ in integer
inverse powers.  At each order the new central value appears
with coefficient one; the other terms have already been fixed
at earlier orders.  Derivatives of the new value first occur
one inverse power later, because odd Gaussian moments vanish.
For each fixed $T$, an auxiliary limit $\lambda\to\infty$ in
the differentiated linear residual identities shows that the
coefficients $\mathfrak p_j$ satisfy exactly the real-phase wave
and gauge transports.  After matching orders below $j$, the
remaining launch discrepancy at order $j$ is therefore in
$\ker K_v$.  Choose the free central kernel value to cancel it,
and extend its jets using $\operatorname{Id}-R_kK$.  This keeps
the compact launch support.  Transport is noncharacteristic at
the entry cylinder, so the matched values and equations match
all required normal jets as well; independent value and normal
conditions are not being imposed on a single free datum.

For clarity, a finite normal-jet budget can be chosen explicitly.
If coefficients through degree $J$ must match to total order $m$
at the last level, assign $m_j=m+2(J-j)$ derivatives to coefficient
$j$. Differentiating its transverse first-order transport at most
$m_j-1$ times requires at most $m_j+1$ derivatives of the preceding
coefficient through the second-order wave forcing. These are
available since $m_{j-1}=m_j+2$. Tangential derivatives of the
matched values and induction on transverse multiplicity then
determine all mixed jets. The differentiated gauge forcing is
only first order in the preceding coefficient and fits the same
budget; Bianchi compatibility adds no independent normal datum.
Increasing the finite $J$ requests further fixed-order coefficient
bounds, without changing the number of derivatives applied to the
physical beam at the already fixed output order.

\emph{Remainders at the actual frequency.}
The matching argument has identified the coefficients and their normal
jets. Its auxiliary-frequency limit alone does not estimate the error
at the experiment's frequency $\lambda=\lambda_T$. For that estimate
we use the following finite bounds.  At each fixed Taylor and transport degree, all
coefficient jets, center-change jets, Gaussian Hessians and their
inverses, and the pairing in Equation~\eqref{bg:contraction-inverse}
cost $e^{o(T)}$.  A Taylor error of transverse order $d$ therefore
gains $\lambda_T^{-d/2}e^{o(T)}$ under Gaussian integration.
Choose the finite degrees backwards from the last required
differentiated residual and matching order.  In the collar one
may first restrict to transverse distances at most
$\lambda_T^{-1/2+\varepsilon_2}$, choosing $\varepsilon_2>0$
small for these finite orders; the complementary Gaussian tail
is negligible.  These finite remainder estimates give
$A_T\lambda_T^{-M_1}e^{o(T)}$ at the actual diagonal frequency.

Realize the finite jets with smooth cutoffs at tube radius
$r_T=e^{-b_pT}$, where $0<2b_p<\zeta$.  On the cutoff annulus,
the higher Taylor terms are smaller than half the positive
quadratic imaginary part for large $T$: their relative bound is
$r_Te^{o(T)}\to0$, using the inverse-Hessian bound as well as
the finite higher-jet bounds.  Positivity therefore implies
\begin{equation}\label{bg:gaussian-tail}
 \exp(-\lambda_T\operatorname{Im}\varphi)
 \le \exp\left[-c\exp\{(\zeta-2b_p)T-o(T)\}\right].
\end{equation}
Every fixed number of cutoff derivatives costs only $e^{Cb_pT}$,
and every fixed phase derivative costs a fixed power of
$\lambda_T$ times $e^{o(T)}$.  Equation~\eqref{bg:gaussian-tail}
dominates all such losses and every prescribed inverse power.
Uniqueness and separation of the entry crossing exclude additional
entry traces.  Tubes may be stopped after they pass through the
strict deeper exit and, on the bridge, lie within arbitrarily
thin margins of Equation~\eqref{bg:arrival-shell}.

Finally absolute integration over the fixed three-dimensional
center patch costs $\lambda_T^{3/2}$; $N_b$ differentiations cost
at most $\lambda_T^{N_b}$ times subexponential factors.  Each
additional finite formal coefficient has subexponential fixed
jets and an assigned nonpositive frequency power.  Increasing
formal accuracy therefore changes constants and the large-$T$
threshold, but neither positive power in the proposition.
Taking real parts, undoing trace reversal, and multiplying by
$A_T$ gives $h_T^B$ and all the claimed bounds.  Every step up to
this point is linear and homogeneous in this common factor.
\end{proof}

\begin{remark}[Order of choices and nonlinear use]\label{bg:choices}
In the application all output, data, trace, constraint, and energy
buffer orders are fixed first.  One then chooses $\zeta>0$ small
enough for their finitely many positive frequency powers, and only
then chooses the finite inverse accuracy and formal degrees.
For example the linear scaled residual can subsequently be made
at most $e^{-4\kappa T}$ by requiring
$\alpha+M_1\zeta>4\kappa$ and then taking $T$ large.
If a nonlinear residual of Equation~\eqref{bg:wave-reduction} is
also evaluated on the exterior beam, Taylor's formula adds a
quadratic term bounded at base order $N_b$ by
$A_T^2\lambda_T^{2N_b+7}e^{o(T)}$, allowing two extra derivatives;
this power again depends only on the fixed base buffer.  The
analogous interior quadratic power is fixed by its base orders.
The signed nonlinear estimates are proved in the following
sections.  The finite tidal observation and deep comparison
conclusions of \cite{Foundation}, which use a different amplitude
and an additional bounded-curvature hypothesis, are not inputs
to Proposition~\ref{bg:packet}.
\end{remark}


\begin{corollary}[Linear construction at any exponentially small amplitude]
\label{bg:amplitude-independent}
Fix any $\alpha>0$ and put $A_T=e^{-\alpha T}$. With this substitution,
all linear conclusions of Proposition~\ref{bg:packet} hold, with the
same order of finite output, frequency and inverse-accuracy choices.
The branch observer needs no bounded-curvature hypothesis.
\end{corollary}
\begin{proof}
Every transport, beam superposition, residual and matching operation
in the proof of Proposition~\ref{bg:packet} is linear and homogeneous
in the final factor $A_T$. Coefficient estimates are independent of
that factor. Multiplying the resulting unit-amplitude construction by
$e^{-\alpha T}$ proves the assertion. In particular, a requested
bound $A_T\lambda_T^{-M_1}e^{o(T)}\le e^{-4\kappa T}$ follows by
choosing the finite accuracy with $\alpha+M_1\zeta>4\kappa$ and then
taking $T$ large.
\end{proof}
